\textbf{Phần 1. Trắc nghiệm lựa chọn}
\setcounterpageref{ex}{0}

\begin{ex} %[4H7V35]
Góc có số đo bằng $90^\circ$ là
\choice
{Góc nhọn}
{\True Góc vuông}
{Góc tù}
{Góc bẹt}
\loigiai{
Góc $90^\circ$ là góc vuông.\\
\textbf{Đáp án của câu hỏi là B}
}
\end{ex}

\begin{ex} %[4H7V35]
Góc bẹt có số đo bằng
\choice
{$90^\circ$}
{$270^\circ$}
{$360^\circ$}
{\True $180^\circ$}
\loigiai{
Góc bẹt $= 180^\circ$.\\
\textbf{Đáp án của câu hỏi là D}
}
\end{ex}

\begin{ex} %[4H7V35]
Hình thoi có
\choice
{\True 4 cạnh bằng nhau và 2 đường chéo vuông góc với nhau}
{4 góc vuông}
{2 đường chéo bằng nhau}
{Chỉ có 1 cặp cạnh song song}
\loigiai{
Hình thoi có 4 cạnh bằng nhau và 2 đường chéo vuông góc.\\
\textbf{Đáp án của câu hỏi là A}
}
\end{ex}

\begin{ex} %[4H7V35]
Một hình bình hành có độ dài đáy $18\text{ cm}$, chiều cao $10\text{ cm}$. Diện tích của nó là
\choice
{$90\text{ cm}^2$}
{$28\text{ cm}^2$}
{\True $180\text{ cm}^2$}
{$56\text{ cm}^2$}
\loigiai{
$S = 18 \times 10 = 180\text{ cm}^2$.\\
\textbf{Đáp án của câu hỏi là C}
}
\end{ex}

\begin{ex} %[4H7V35]
Chu vi hình thoi có cạnh dài $15\text{ cm}$ là
\choice
{$30\text{ cm}$}
{\True $60\text{ cm}$}
{$45\text{ cm}$}
{$225\text{ cm}$}
\loigiai{
$P = 15 \times 4 = 60\text{ cm}$.\\
\textbf{Đáp án của câu hỏi là B}
}
\end{ex}

\begin{ex} %[4H7V35]
Hai đường thẳng song song là hai đường thẳng
\choice
{Cắt nhau tại góc $90^\circ$}
{Cắt nhau tại 1 điểm}
{Trùng nhau}
{\True Không có điểm chung nào (dù kéo dài mãi)}
\loigiai{
Hai đường thẳng song song không cắt nhau.\\
\textbf{Đáp án của câu hỏi là D}
}
\end{ex}

\begin{ex} %[4H7V35]
Lúc 6 giờ đúng, hai kim giờ và kim phút tạo thành góc có số đo là
\choice
{\True $180^\circ$}
{$90^\circ$}
{$120^\circ$}
{$60^\circ$}
\loigiai{
Góc bẹt $180^\circ$.\\
\textbf{Đáp án của câu hỏi là A}
}
\end{ex}

\begin{ex} %[4H7V35]
Một mảnh đất hình bình hành có đáy $30\text{ m}$, chiều cao bằng một nửa đáy. Diện tích mảnh đất là
\choice
{$225\text{ m}^2$}
{$900\text{ m}^2$}
{\True $450\text{ m}^2$}
{$300\text{ m}^2$}
\loigiai{
Chiều cao: $30 : 2 = 15\text{ m}$.\\
Diện tích: $30 \times 15 = 450\text{ m}^2$.\\
\textbf{Đáp án của câu hỏi là C}
}
\end{ex}

\begin{ex} %[4H7V35]
Một hình thoi có chu vi bằng $72\text{ dm}$. Độ dài một cạnh hình thoi là
\choice
{$16\text{ dm}$}
{\True $18\text{ dm}$}
{$24\text{ dm}$}
{$36\text{ dm}$}
\loigiai{
Cạnh: $72 : 4 = 18\text{ dm}$.\\
\textbf{Đáp án của câu hỏi là B}
}
\end{ex}

\begin{ex} %[4H7V35]
Trong tam giác $ABC$ có góc $A = 90^\circ$, góc $B = 35^\circ$. Số đo góc $C$ là
\choice
{$45^\circ$}
{$65^\circ$}
{$50^\circ$}
{\True $55^\circ$}
\loigiai{
Góc $C = 90^\circ - 35^\circ = 55^\circ$.\\
\textbf{Đáp án của câu hỏi là D}
}
\end{ex}

\begin{ex} %[4H7V35]
Số cặp cạnh đối diện song song trong một hình chữ nhật là
\choice
{\True $2$ cặp}
{$1$ cặp}
{$3$ cặp}
{$4$ cặp}
\loigiai{
Có 2 cặp cạnh đối diện song song.\\
\textbf{Đáp án của câu hỏi là A}
}
\end{ex}

\begin{ex} %[4H7V35]
Một mảnh vườn hình chữ nhật có chu vi $100\text{ m}$, chiều dài hơn chiều rộng $10\text{ m}$. Người ta xây một lối đi hình bình hành chạy xuyên qua vườn có đáy $2\text{ m}$ và chiều cao bằng chiều rộng mảnh vườn. Diện tích lối đi là
\choice
{$30\text{ m}^2$}
{$50\text{ m}^2$}
{\True $40\text{ m}^2$}
{$60\text{ m}^2$}
\loigiai{
Nửa chu vi: $50\text{ m}$. Dài: $30\text{ m}$, Rộng: $20\text{ m}$.\\
Diện tích lối đi: $2 \times 20 = 40\text{ m}^2$.\\
\textbf{Đáp án của câu hỏi là C}
}
\end{ex}

\textbf{Phần 2. Trắc nghiệm đúng sai}
\setcounterpageref{ex}{0}

\begin{ex} %[4H7V35]
Xét tính đúng sai về các hình học đã học:
\choiceTF
{\True Góc nhọn $<$ Góc vuông $<$ Góc tù $<$ Góc bẹt.}
{\True Hai đường chéo hình thoi vuông góc với nhau.}
{Hình bình hành luôn có 4 cạnh bằng nhau.}
{\True Hình vuông vừa là hình chữ nhật đặc biệt vừa là hình thoi đặc biệt.}
\loigiai{
\begin{itemize}
    \item Thứ tự góc: Đúng.
    \item Hai đường chéo hình thoi vuông góc: Đúng.
    \item Hình bình hành chỉ có các cạnh đối bằng nhau. Sai.
    \item Hình vuông là hình chữ nhật và hình thoi đặc biệt: Đúng.
\end{itemize}
\begin{itemize}
    \item a) \textbf{Đúng}.
    \item b) \textbf{Đúng}.
    \item c) \textbf{Sai}.
    \item d) \textbf{Đúng}.
\end{itemize}
}
\end{ex}

\begin{ex} %[4H7V35]
Cho hình bình hành $MNPQ$ có đáy $PQ = 20\text{ cm}$, chiều cao tương ứng $MH = 12\text{ cm}$. Xét tính đúng sai:
\choiceTF
{\True Đoạn thẳng $MH$ vuông góc với cạnh đáy $PQ$.}
{\True Cạnh $MN$ song song và bằng cạnh $PQ = 20\text{ cm}$.}
{Diện tích hình bình hành là $120\text{ cm}^2$.}
{\True Diện tích hình bình hành là $240\text{ cm}^2$.}
\loigiai{
\begin{itemize}
    \item $MH \perp PQ$: Đúng.
    \item $MN // PQ$ và $MN = PQ = 20\text{ cm}$: Đúng.
    \item $S = 20 \times 12 = 240\text{ cm}^2 \neq 120\text{ cm}^2$. Ý 3 Sai, Ý 4 Đúng.
\end{itemize}
\begin{itemize}
    \item a) \textbf{Đúng}.
    \item b) \textbf{Đúng}.
    \item c) \textbf{Sai}.
    \item d) \textbf{Đúng}.
\end{itemize}
}
\end{ex}

\textbf{Phần 3. Trả lời ngắn}
\setcounterpageref{ex}{0}

\begin{ex} %[4H7V35]
Số đo của góc vuông là bao nhiêu độ?
\shortans{90}
\loigiai{
\textbf{Đáp số:} $90$\\
$90^\circ$.
}
\end{ex}

\begin{ex} %[4H7V35]
Tính chu vi hình thoi có cạnh dài $25\text{ cm}$.
\shortans{100}
\loigiai{
\textbf{Đáp số:} $100$\\
$25 \times 4 = 100\text{ cm}$.
}
\end{ex}

\begin{ex} %[4H7V35]
Một hình bình hành có đáy $15\text{ cm}$, diện tích $150\text{ cm}^2$. Tính chiều cao tương ứng theo xăng-ti-mét.
\shortans{10}
\loigiai{
\textbf{Đáp số:} $10$\\
$150 : 15 = 10\text{ cm}$.
}
\end{ex}

\begin{ex} %[4H7V35]
Lúc 4 giờ đúng, góc nhỏ tạo bởi hai kim đồng hồ có số đo bao nhiêu độ?
\shortans{120}
\loigiai{
\textbf{Đáp số:} $120$\\
$4 \times 30^\circ = 120^\circ$.
}
\end{ex}

\begin{ex} %[4H7V35]
Một hình thoi có chu vi $64\text{ cm}$. Tính độ dài cạnh hình thoi theo xăng-ti-mét.
\shortans{16}
\loigiai{
\textbf{Đáp số:} $16$\\
$64 : 4 = 16\text{ cm}$.
}
\end{ex}

\begin{ex} %[4H7V35]
Một hình vuông có cạnh $10\text{ cm}$. Người ta nối trung điểm các cạnh của hình vuông đó để được một hình thoi bên trong. Tính diện tích hình thoi đó theo xăng-ti-mét vuông.
\shortans{50}
\loigiai{
\textbf{Đáp số:} $50$\\
Diện tích hình thoi bằng một nửa diện tích hình vuông: $10 \times 10 : 2 = 50\text{ cm}^2$.
}
\end{ex}
